\begin{abstract}
\noindent
楔形变截面矩形拉杆的轴向刚度与应力分布，是检验一维杆理论适用边界、以及二维单元离散质量的典型问题。本文以截面高度线性变化、关于杆轴对称的矩形拉杆为对象，取 $\beta=0$ 的等截面情形作为退化校核、$\beta=+1.0$ 为主构型，分别实现一维变截面杆单元、四节点双线性等参元 Q4 和三节点常应变三角形 CST，并用 CalculiX 的 CPS4/CPS8 平面应力模型与 FEniCS/dolfinx 作第三方锚点。

数值结果表明：在 $160\times32$ 网格上，Q4 的端部、中截面平均位移相对 CalculiX CPS8（$240\times48$）参考的最大偏差分别为 $0.025\%$ 和 $0.014\%$，探针点应力偏差不超过 $0.042\%$；CST 对应为 $0.075\%$、$0.067\%$ 和 $0.146\%$。Q4 的梯度保真误差为 $0.35\%$–$0.95\%$，位移与应力的观测收敛阶分别约为 $1.7$–$2.5$ 和 $1.6$–$2.3$。CST 的两种对角线剖分在对称工况下互为镜像，端/中位移与截面平均应力是严格不变量（离散度约 $10^{-11}\%$），真正的对角线敏感性只出现在局部单元/节点应力（$160\times32$ 下为 $0.2\%$–$1.1\%$）。

一维杆判据采用 $\eta=|\beta|h_0/L=\tan\alpha$，将工况分为 A（可信）、B（需二维复核）、C（不可信）三级；与 CalculiX 参考相比，一维杆端部偏差依次为 $+0.25\%$、$-0.55\%$、$+0.05\%$、$-3.21\%$、$-7.36\%$，与标签分级一致。CalculiX 同网格 CPS4 与自研 Q4 的端位移差不超过 $0.05\%$；FEniCS 对 S1 的端位移与 CCX 参考仅差 $+0.0029\%$。三种方法在产出形态上的差异——一维曲线、连续云图、分块云图——与它们的精度、自由度和后处理需求同样重要，本文为此单列一节对比。
\end{abstract}

\noindent\textbf{关键词：}变截面杆；平面应力；Q4 等参元；CST；等效刚度；收敛性；CalculiX；对角线敏感性
